\chapter{The Insufficiency of Flat Maps}

What became clear to me, as I pondered this question with the seriousness it deserves, is that the raging debate about the actual shape of the world—flat or spherical or some other exotic geometry—has entirely missed the deeper point. All our current map projections, every single one of them, fail to contain the information we truly need. They are equally insufficient in their capacity to describe what is really there; they can only describe what it looks like from a certain vantage, and often a rather limited one at that. They capture the light photons bouncing off surfaces, but not the deeper content beneath.

In addition—and this is perhaps the more troubling observation—they fail to grow and develop with humanity. They have scarcely changed in centuries, really, yet we continue to treat them as though they represent the state of the art. We do not evolve our knowledge of mapping the world around us; instead, we seem compelled to take these ancient maps and force our growing, gargantuan knowledge of the world to conform to them. That is a strange and almost tragic inversion.

So I thought: if we take the shape of the knowledge we actually possess—which is vast and intense, composed of enormous data sets of geo-positions and points of interest and historical boundaries and ownership claims, and that is only the beginning—there is so much there. And that is not even touching the actual physics, the content of the rock itself, how deep it goes, what lies beneath us, what lies above us. There is so much missing from our maps that it is frankly insane we still use them. And even further, it is insane that we would continue to argue between a flat Earth and a spherical Earth when neither model contains any genuinely useful information at all. We have to take all this rich, complex information and cram it into a little tiny shape in order to communicate it or store it. That, I believe, is a problem that can and must be solved.

What became clear when I looked at the maps—and I looked a great many of them—was not that they all contained similar information, but that they all contained similar blocks. They all negated certain understandings; they all pushed you toward a very limited way of looking at the surface. They literally made a surface that was superficial.

And that superficiality pushed itself into more places than we might imagine. It pushed into our games as children, when we drew maps and wandered in and out of the world. It pushed into the games we play online now—Minecraft, Roblox, all of them—each one perfectly reproducing the view of this real world we inhabit. The shape and contour of the world is, paradoxically, made harder to understand by looking at the map.

The question, then, becomes: What is the optimal perspective on information and data—not merely geospatial, but spatial in the broadest sense? If we were to have a truly representational map of everything, what would it be? And the answer, it seemed to me, was obvious: a computer model.

When I examined the models—and I examined a great many of them—I noticed something striking. Almost all of them were based upon correlation of data over time. Because time series were their most common input, the implicit assumption was that a projection into the future could be made that would be actionable. And it is not true. A projection based upon time records and the information the model was fed is useful, perhaps, but it is not actionable. What is needed is a causal model—a model that can accurately draw from real-time information about supply, mass, density, and the like. A truly causal model.

And there are no causal models for local environments. There are no models, really, for local environments that are even correlative; they are all global simulations that run on time-based data and project forward, or they train on other data and produce something vaguely similar to causation, but not quite causal—not granular enough to know at a precise level where everything is.

All the things in our lives are in a place where we know they are. We have an associative map in our minds of where one thing is and how it relates to an activity, how it relates to a person, how it relates to sentiment over the item. We have maps of everything we are aware of, and these are causal maps. These are maps where we have witnessed the inventory, where we know, relative to other things, where other things are—often able to reach out and grasp something without even looking in environments where we are deeply familiar with the objects. I began to wonder: How could we make such a model, the one we already possess in our minds, but for the whole local area around us—and for local areas we do not inhabit?

And this is confirmed by neurological research—into things like Neuralink, for instance—and by our long-established body of research into cognition: channel capacity, the complexity of an object and how it determines how well we can remember it.

So the task in finding the perfect shape for the map to follow is not in finding a shape out in the world, but in making the connections necessary within the information we already have—connections we have attempted to conform to the old flat-map models. The point I was trying to make is that we must look at what we did to the data to force it to fit the flat map model and then flip it inside out, because the flat map itself is already a model of how we mapped our associative maps onto the world. It is a transformer.

What we discover, when we do that, is a big-data problem: too much data and too little display, too small a window for too much.

Now, if we turn back to the sky—that big blue giant above us, with all those stars in it—the same sky the ancient people gazed at, the same sky that means slightly different things to us now, but you get the point. It is big and old and has been largely the same for an extraordinarily long time. Slight changes. Variations.

The sky is a very old thing. It has been above us for quite some time, largely unchanged throughout the entire span of human existence on Earth. You can go all the way back to the oxygenation event and find that once an oxygen-rich atmosphere and biological life were established, there was never any interruption in the capacity to maintain life.

Similarly, there are other things on Earth that are as continuous, that appear normally on a map but whose true entity you would not understand or see. For example, the surface area of the planet maintains itself as a constant. It does not grow or shrink very much at all; it has not in millennia. Yet we can observe seafloor spreading; we can observe subduction.

There are plenty of surface-area maps, but none of them will tell you the magical truth that this surface area is maintained somehow as a constant while it is constantly changing. And then I learned a new term: dynamic steady-state equilibrium—a concept that, once grasped, reshapes how we see not just the Earth’s crust, but the very structure of information itself.

On learning what a steady-state dynamic equilibrium was, I decided to take another look at the maps—to see if I could find such an equilibrium on the map itself. And I found a quite profound steady-state dynamic equilibrium on our written flat paper maps.

I found it in all the maps we have ever written. It is the steady-state dynamic equilibrium where the steady-state is the human need to describe the area around us, and the dynamic equilibrium is the information that is added and removed from different maps produced by different cultures and for different purposes. There is always a map of the world, and that map is always taking a particular shape, and that map is used for many different things, and information goes into it and comes out of it.

So, it is not really that the map itself is a state-of-the-art tool that needs upgrading. The map explains something about people—something about us—something about our own steady-state dynamic equilibrium of information.

So, in seeing the map differently now—seeing it not merely as the map itself but as a revelation of how humans are with maps, what maps truly are to us—I thought that I would look again at the sky. I thought I would examine sky maps, to see how we have depicted the heavens above us, to determine if there was some correlation between the way we have rendered the terrestrial surface and the way we have rendered the celestial. And of course there was—but it was odd, profoundly odd. It involved a great deal of superstition, a pantheon of gods, old ideas long discarded (or so we like to think). A great many pantheons, in fact, all infused with superstition.

There is, I came to realize, a steady-state dynamic equilibrium in our maps of the sky—not one where the steady-state is a particular shape or territory, but where the steady-state is the presence of the infinite unknown. We are all aware, at some deep level, that there is far more up there than we can grasp, and so our maps tend toward constellations, toward patterns imposed upon the void. They do not tend toward defined territories or bounded areas, because we know, instinctively perhaps, that the sky extends to infinity in one direction—away.

Sky maps are, therefore, particularly prone to superstition: constellations named after gods, planets named after gods (and let us remember that “planet” originally meant “wanderer”—those lights that moved against the fixed backdrop of the stars).

So now we can see the sky as a dynamic steady-state equilibrium of information about the cosmos—one deliberately infused with superstition and religion, precisely because it is the plane upon which we confront the infinite. And since we cannot fully comprehend or contain the infinite, there is necessarily an element of faith involved—an element of trust. A person with a telescope can see farther than the rest of us, and we approach that disparity with no small amount of fear within the orthodoxy. And rightly so—because if you can see farther than everyone else, you possess power.

We are, in this sense, always looking at people looking at each other through their maps. On the surface of the Earth, we have constructed a steady-state dynamic equilibrium of the information we deem necessary or important at the time—not all the information we possess, by any means, but the information vital for navigation during eras of expansion and exploration, and now for things like agriculture, deforestation, carbon tracking as it moves through the atmosphere. We have many different ways of using surface maps now, and they are generally not imbued with religious connotation.

The sky, however, is a different matter entirely—particularly now that we possess large space-based telescopes capable of observing from beyond the distorting veil of the atmosphere. The distance from the average person to the information yielded by these instruments is vast. There is much we must trust these “seers in the sky” to reveal on our behalf. And so a large portion of modern science—stellar physics, astrophysics—has taken on a kind of reverence, a religious aura, as if it has inherited the mantle of the sacred.

This is perhaps best illustrated by the extraordinary images from the Hubble and now the James Webb Telescope—beautiful, almost transcendent mosaics of scenes unimaginably far away. The average person gazes at them in awe but has little idea of the perspective: Where are we in the photograph? Where do we fit into the picture? That crucial context is missing. And between us and that image stands a medium—a mediator, a middleman.

So let us examine the middleman. Where does he reside? Let us return, for a moment, to the Apollo moon landings—to the photograph of Earth taken from the lunar surface.

Earthrise was a milestone in human achievement, a genuine milestone. From that image, taken from the Moon, you can see the Earth rising above the barren horizon; you can see where we fit into the picture.

At that point in history, when the Apollo photographs emerged, it became possible—for the first time, perhaps—to truly see where we fit into the universe. And the mediator, the middleman, was an astronaut: a physical human being on the Moon, someone we knew was there because, look, there is the picture of us taken from the Moon.

Earthrise was a moment when we could still see ourselves in the picture. You could see the astronauts, you could see our fragile blue Earth, and you could see where we fit into the broader cosmos. Yet we still had the tangible element of the astronaut and the infinite backdrop of the sky.

It is therefore no surprise that, one or two generations later, an amount of doubt and skepticism arose regarding something that had occurred, that carried scientific fact, and that then gave rise to a kind of orthodoxy—one that would swear allegiance without need for further evidence and reject any questioning of the accepted truth. Not because scientific rigor had been lost, and not because counter-arguments possessed superior facts, but simply because the discussion touched upon the dynamic equilibrium that includes infinity—the realm we map not with science alone, for we are deeply prone to regarding the sky with superstition.

Venturing to the Moon provided powerful confirmation—perhaps even confirmation bias—that the information we store in the sky, or the relatives we imagine there, is prone to battles between non-scientific forces of skepticism and belief.

These debates—about events considered scientific fact that devolve into orthodoxy of belief versus skepticism seeking evidence that does not exist or has already been found—are not unique to any particular controversy. They are simply a feature of information stored within the dynamic equilibrium of mapping the sky.

In addition, science provides us with a continuously changing landscape of information—new discoveries constantly emerging, old certainties constantly overturned. We possess a truly dynamic information set regarding what we see in the sky.

The more mundane, the earthly reason for the separation between these two dynamics is that the information we place on maps of the Earth and our immediate surroundings is directly tied to our lives. It influences the shape, the condition, the quality of our existence in tangible ways.

In a sense, the things we can view and empirically experience in our environment are mapped differently from the things we must experience abstractly in the sky above us. Direct contact demands one kind of mapping; the infinite demands another. But what about the things beneath the Earth? What about the interior?

Let us consider how far we have penetrated the crust.

Seven point six miles—approximately. It is a strikingly small figure; one might think we would have gone farther. Yet the maximum depth we have achieved is roughly 7.6 miles.

The distance from the surface to the center of the Earth is, by contrast, nearly 4,000 miles. The full diameter, crust to antipodal crust, is almost 8,000 miles.

We have penetrated less than one-fifth of one percent toward the center—less than one-tenth of one percent all the way through.

The average distance to the Moon is approximately 239,000 miles. The International Space Station orbits at about 250 miles above the surface. We have ventured vastly farther upward than downward.

With only two geo-positions of note—the Challenger Deep at roughly 6.8 miles below sea level and the Kola borehole at 7.6 miles below the continental surface—we possess merely two points of direct exploration into a planet full of unseen, unmapped substance beneath our feet.

So let us return, with all this knowledge in hand, to our dynamic steady-state equilibrium. We now understand that there is a crustal surface-area dynamic steady-state equilibrium, and we understand that there is a dynamic steady-state equilibrium in the atmosphere.

Let us see if we can use our maps of these equilibria to project inward—to discern what might lie beneath the crustal surface-area equilibrium maintained by subduction and spreading rifts.

To see what is beneath, we need only examine where this vast, uninterrupted steady state has been disrupted by events.

And we arrive at a construct: the volcano.

What occurs in a volcanic eruption is profoundly interesting when viewed in this light. It affects one of these large, uninterrupted steady-state dynamic equilibria that have persisted throughout geologic time. How does the system maintain its balance when confronted with a massive plume of activity—rock dust, lava, gas, material surging forth? What is truly happening?

There is a visible clue: ejecta in the form of lava shot outward, landing on the surface—a surface we know will eventually subduct back into the mantle over vast timescales.

There is, therefore, a dynamic entity beneath the crust capable of depositing material upon the surface, material that will ride the conveyor of plate tectonics back down into the depths.

We can see that beneath the crustal steady-state dynamic equilibrium of surface area—where subduction and spreading rifts operate—something breaks through the crust, deposits energy and matter on the surface, matter and energy that will eventually return to the mantle. Some deeper steady-state is being maintained through the management of dynamic entities like the volcano—a transfer of energy from one steady-state dynamic equilibrium to another. These are, in essence, parallel or nested equilibria.

The final observation—and it is a crucial one—is that the volcanic eruption represents what is known as entropic heat loss.

But not merely entropic heat loss. There is entropic process matched with negentropic deposition: the eruption dissipates heat, yet it deposits structured material—everything transferred from one layer to the next, material destined to return. And in the sky, when entropic heat loss occurs (evaporation, dispersal), condensation follows, and rain falls—negentropic deposition upon the surface.

The profound similarity between these transitions lies in the one-to-one energy relationship between the layers conducting the transfer. The volcano and the rain are, in a deep sense, the same phenomenon: transitions between distinct steady-state dynamic equilibria where energy is moved from one system to another, time-shifted, destined to return through subtler means. These one-to-one transfers are, it seems to me, a fundamental thing.

And what that thing is—is a steady-state dynamic equilibrium paradigm that transfers energy between itself and similar entities. It is older than everything else, uninterrupted, capable of producing all the wonders we see. Mapping this thing would allow us—would allow you and me—to map the entirety of the world.

I would like to conclude this reflection with the idea that this nested steady-state dynamic equilibrium model stands as the optimal format for dynamic data in causal model-capable data storage. It is the shape—the container—that can hold all our vast, ever-shifting knowledge while preserving the causal connections, the associative depths, that flat maps or mere correlative models so woefully distort. This paradigm, with its layered transfers and balanced equilibria, apparently contains everything on our planet. Everything on our planet is encompassed by such a structure, one that accommodates this type of transfer and this type of steady-state dynamic equilibrium—making it not just a map, but the very architecture for relating our lives to the real.

And the final step—for our understanding, perhaps for the cosmos itself—is this: once we possess a map of that thing, we can decide what to call it. We can then project that map outward into the infinity beyond us in the sky, and we can begin to see these paradigms as the shape of planets everywhere—as each planet a set of nested equilibria conforming to this model. This thing that, if we truly map it, allows us to map everything else.
